% !PRINTED 189
\weberpage{189}

Let it finally be remarked that, instead of the \gesperrt{positions} of the
body, the rotations themselves can also be conceived as elements
of the group.

\begin{center}
\S.\,49.

\gesperrt{Linear fractional substitutions.}
\end{center}

The group of orthogonal ternary substitutions is, as shall
now be proved, isomorphic with the group of linear fractional
substitutions of one variable, or of the binary substitutions of
the ratios (\S.\,38).

We denote a ternary orthogonal substitution
by
\weq{1}{(y_1,\ y_2,\ y_3) = A\,(x_1,\ x_2,\ x_3),}
\weq{2}{y_1^2 + y_2^2 + y_3^2 = x_1^2 + x_2^2 + x_3^2.}
To this we first apply, following \S.\,37, (21), the
transformation by the fixed (not orthogonal) substitution
\weq{3}{L = \left\{\begin{array}{ccc}
i, & 0, & 0 \\[4pt]
0, & \dfrac{1}{\sqrt{2}}, & \dfrac{1}{\sqrt{2}} \\[6pt]
0, & \dfrac{i}{\sqrt{2}}, & \dfrac{-i}{\sqrt{2}}
\end{array}\right\},\quad
L^{-1} = \left\{\begin{array}{ccc}
-i, & 0, & 0 \\[4pt]
0, & \dfrac{1}{\sqrt{2}}, & \dfrac{-i}{\sqrt{2}} \\[6pt]
0, & \dfrac{1}{\sqrt{2}}, & \dfrac{i}{\sqrt{2}}
\end{array}\right\}}
in which $i = \sqrt{-1}$, and whose determinant has the value 1,
and set
\weq{4}{\begin{aligned}
(y_1,\ y_2,\ y_3) &= L\,(y_1',\ y_2',\ y_3')\\
(x_1,\ x_2,\ x_3) &= L\,(x_1',\ x_2',\ x_3')' % [sic] source prints a stray prime after the closing parenthesis
\end{aligned}}
\weq{5}{(y_1',\ y_2',\ y_3') = A'\,(x_1',\ x_2',\ x_3'),}
in which
\weq{6}{A' = L^{-1} A L}
likewise denotes a linear substitution, whose determinant
$|A'| = |A|$, and hence is equal to $\pm 1$. The relation
(2) passes, by the substitutions (4), into
\weq{7}{-y_1'^2 + 2\,y_2'\,y_3' = -x_1'^2 + 2\,x_2'\,x_3'}
from which six relations between the coefficients of $A'$
result, which however we shall not set up here, since we can
find the substitution $A'$ more simply in another way.
